\chapter{LLD — Orchestration des quotas}
\label{chap:orchestration}

\section{Objet}
Ce chapitre spécifie le modèle de transfert réversible de quota individuel vers un projet de groupe, retenu comme modèle cible de gouvernance des ressources pour les étudiants.

\section{Modèle retenu — Modèle C}
\begin{infobox}[Modèle de transfert de quota — Modèle C]
\begin{enumerate}
  \item Chaque étudiant dispose d'un projet Keystone individuel avec un quota standard.
  \item Un groupe est un projet Keystone séparé, dont le quota est la somme dynamique des contributions de ses membres.
  \item Un étudiant contribue une partie de son quota \textbf{libre} à un ou plusieurs groupes, dans la limite de son quota individuel.
  \item La contribution est réversible : restituée lorsque le groupe repasse sous la nouvelle limite, ou lors de la suppression du groupe.
  \item Tant que le groupe reste au-dessus de la nouvelle limite, la contribution est \textbf{gelée}.
\end{enumerate}
\end{infobox}

\section{Règle de gel — anti-abus}
Soit :
\begin{itemize}
  \item $L$ : ancienne limite du groupe (somme des contributions actives).
  \item $c$ : contribution du partant.
  \item $L' = L - c$ : nouvelle limite.
  \item $U$ : usage courant du groupe (vCPU effectivement consommés).
\end{itemize}

\begin{tabularx}{\textwidth}{@{}lXX@{}}
\toprule
\textbf{Condition} & \textbf{État contribution} & \textbf{Quota étudiant} \\
\midrule
$U \leq L'$ & \texttt{RELEASED} immédiatement & Restitué \\
$U > L'$    & \texttt{RELEASE\_PENDING} & Gelé (non restitué) \\
\bottomrule
\end{tabularx}

La libération n'est jamais partielle : soit les $c$ vCPU sont entièrement restitués, soit ils restent entièrement gelés.

\section{Modèle de données}
\begin{lstlisting}[language=SQL, caption={Schéma PostgreSQL}]
CREATE TABLE student (
  id                  UUID PRIMARY KEY,
  login               VARCHAR(64)  UNIQUE NOT NULL,
  email               VARCHAR(255),
  keystone_user_id    UUID NOT NULL,
  keystone_project_id UUID NOT NULL,
  quota_vcpu          INTEGER NOT NULL DEFAULT 4,
  quota_ram_mb        INTEGER NOT NULL DEFAULT 8192,
  quota_disk_gb       INTEGER NOT NULL DEFAULT 100,
  state               VARCHAR(16) NOT NULL DEFAULT 'ACTIVE',
  created_at          TIMESTAMPTZ NOT NULL DEFAULT now()
);

CREATE TABLE grp (
  id                  UUID PRIMARY KEY,
  name                VARCHAR(128) NOT NULL,
  keystone_project_id UUID NOT NULL,
  state               VARCHAR(16) NOT NULL DEFAULT 'ACTIVE',
  created_at          TIMESTAMPTZ NOT NULL DEFAULT now(),
  pending_delete_at   TIMESTAMPTZ,
  deleted_at          TIMESTAMPTZ
);

CREATE TABLE contribution (
  id                   UUID PRIMARY KEY,
  student_id           UUID NOT NULL REFERENCES student(id),
  group_id             UUID NOT NULL REFERENCES grp(id),
  vcpu                 INTEGER NOT NULL CHECK (vcpu > 0),
  ram_mb               INTEGER NOT NULL CHECK (ram_mb > 0),
  disk_gb              INTEGER NOT NULL CHECK (disk_gb > 0),
  state                VARCHAR(24) NOT NULL DEFAULT 'COMMITTED',
  created_at           TIMESTAMPTZ NOT NULL DEFAULT now(),
  release_requested_at TIMESTAMPTZ,
  released_at          TIMESTAMPTZ,
  release_reason       VARCHAR(64)
);

CREATE TABLE audit_log (
  id            BIGSERIAL PRIMARY KEY,
  ts            TIMESTAMPTZ NOT NULL DEFAULT now(),
  actor         VARCHAR(128),
  action        VARCHAR(64) NOT NULL,
  entity_type   VARCHAR(32) NOT NULL,
  entity_id     UUID NOT NULL,
  before_state  JSONB,
  after_state   JSONB,
  result        VARCHAR(16) NOT NULL,
  error_message TEXT
);
\end{lstlisting}

\section{Machines à états}
\begin{figure}[htbp]
\centering
% Machine à états d'une contribution
\begin{tikzpicture}[
  st/.style={base, rectangle, minimum width=3.6cm, minimum height=1.3cm, font=\ttfamily\footnotesize\bfseries},
]
  \node[st, fill=accent!8, line width=1.1pt] (c) at (0,0)    {COMMITTED\\[1pt]\sffamily\mdseries\scriptsize membre actif};
  \node[st]                 (p) at (8.4,0)  {RELEASE\_PENDING\\[1pt]\sffamily\mdseries\scriptsize quota gelé};
  \node[st]                 (r) at (4.2,-3.8) {RELEASED\\[1pt]\sffamily\mdseries\scriptsize quota restitué};

  \draw[flow,->] ([yshift=7pt]c.east) -- node[lab, above=1pt] {départ, $U > L'$} ([yshift=7pt]p.west);
  \draw[flow,->] ([yshift=-7pt]p.west) -- node[lab, below=1pt] {réintégration} ([yshift=-7pt]c.east);
  \draw[flow,->] (c.south) |- node[lab, above=1pt, pos=0.75] {départ, $U \leq L'$} (r.west);
  \draw[flow,->] (p.south) |- node[lab, above=1pt, pos=0.75, align=center] {$U \leq L'$ ou\\groupe supprimé} (r.east);
\end{tikzpicture}

\caption{Machine à états d'une contribution.}
\label{fig:etat-contribution}
\end{figure}

\begin{figure}[htbp]
\centering
% Machine à états d'un groupe
\begin{tikzpicture}[
  st/.style={base, rectangle, minimum width=3.2cm, minimum height=1.3cm, font=\ttfamily\footnotesize\bfseries},
]
  \node[st, fill=accent!8, line width=1.1pt] (a) at (0,0)     {ACTIVE\\[1pt]\sffamily\mdseries\scriptsize $\geq 2$ membres};
  \node[st]                 (p) at (5.7,0)   {PENDING\_DELETE\\[1pt]\sffamily\mdseries\scriptsize $< 2$ membres, 30 j};
  \node[st]                 (d) at (11.4,0)  {DELETED\\[1pt]\sffamily\mdseries\scriptsize projet supprimé};

  \draw[flow,->] ([yshift=7pt]a.east) -- node[lab, above=1pt] {effectif $< 2$} ([yshift=7pt]p.west);
  \draw[flow,->] ([yshift=-7pt]p.west) -- node[lab, below=1pt] {nouveau membre} ([yshift=-7pt]a.east);
  \draw[flow,->] (p.east) -- node[lab, above=1pt, align=center] {délai expiré\\ou action admin} (d.west);
\end{tikzpicture}

\caption{Machine à états d'un groupe.}
\label{fig:etat-groupe}
\end{figure}

\section{Invariants}
\begin{enumerate}
  \item \textbf{Étudiant} — $\sum$ contributions actives $\leq$ quota individuel.
  \item \textbf{Contribution} — On ne contribue que du quota libre : usage + contributions + nouvelle $\leq$ quota.
  \item \textbf{Groupe} — Le quota du groupe est exactement la somme de ses contributions actives.
  \item \textbf{OpenStack} — Les quotas appliqués à chaque mutation sont recalculés depuis la base.
  \item \textbf{Restitution} — Une contribution en \texttt{RELEASE\_PENDING} bascule en \texttt{RELEASED} en bloc, jamais partiellement.
\end{enumerate}

\begin{infobox}[Nature du quota transféré]
Le quota transféré est un quota \textbf{logique}, pas une réservation physique. OpenStack gère l'overcommit via ses ratios. La contribution transfère un droit de créer des instances, pas des cœurs physiques.
\end{infobox}

\section{API d'orchestration}
\begin{tabularx}{\textwidth}{@{}lllX@{}}
\toprule
\textbf{Méthode} & \textbf{Chemin} & \textbf{Rôle} & \textbf{Description} \\
\midrule
POST   & /students                                 & admin        & Créer étudiant \\
GET    & /students/\{id\}/available                & self / admin & Quota libre contribuable \\
POST   & /groups                                   & student      & Créer groupe + contribution \\
POST   & /groups/\{id\}/contributions              & student      & Rejoindre un groupe \\
DELETE & /groups/\{id\}/contributions/\{cid\}      & self / admin & Quitter le groupe \\
POST   & /groups/\{id\}/contributions/\{cid\}/rejoin & self       & Réintégration \\
GET    & /groups/\{id\}/usage                      & member / admin & Usage et quota du groupe \\
POST   & /admin/groups/\{id\}/force-release/\{cid\} & admin       & Libération forcée (escalade) \\
GET    & /audit                                    & admin        & Journal d'audit \\
\bottomrule
\end{tabularx}

\section{Jobs de fond}
\begin{tabularx}{\textwidth}{@{}lXl@{}}
\toprule
\textbf{Job} & \textbf{Rôle} & \textbf{Fréquence} \\
\midrule
reconcile\_quota\_drift    & Compare DB et OpenStack, corrige les écarts & 15 min \\
check\_release\_pending    & Bascule en \texttt{RELEASED} si $U \leq L'$ & 5 min \\
check\_group\_dissolution  & Gère le délai 30 j des \texttt{PENDING\_DELETE} & 1 h \\
notify\_pending\_release   & Notifications J+0, J+7, J+14 & 1×/jour \\
expire\_groups             & Nettoyage et purge des contributions anciennes & 1×/jour \\
\bottomrule
\end{tabularx}

\section{Cohérence et réconciliation}
Pattern outbox : la transaction DB écrit la contribution et un enregistrement outbox. Un worker consomme l'outbox et applique l'action OpenStack. En cas d'échec : retry avec backoff (5 max), puis alerte. Job périodique de réconciliation pour les écarts résiduels.

\begin{okbox}[Garantie]
Aucune écriture OpenStack sans enregistrement outbox préalable. Aucune contribution considérée appliquée tant que l'outbox n'est pas \texttt{DONE}.
\end{okbox}

\section{Sécurité et déploiement}
\begin{itemize}
  \item Compte de service OpenStack stocké dans un vault.
  \item API exposée uniquement sur le management, derrière SRX + VPN.
  \item Rate limiting : 10 req/min par étudiant sur les contributions.
  \item Journal d'audit exhaustif.
\end{itemize}

\begin{tabularx}{\textwidth}{@{}lrrrX@{}}
\toprule
\textbf{Composant} & \textbf{vCPU} & \textbf{RAM} & \textbf{Disque} & \textbf{Remarque} \\
\midrule
VM management (API + worker + Redis) & 2 & 4 Go & 40 Go & Phase pilote \\
PostgreSQL                            & 2 & 4 Go & 80 Go & Séparer en production \\
Worker Celery                         & 1 & 2 Go & 20 Go & Cohabite en pilote \\
\bottomrule
\end{tabularx}